\section{L7 -- CPU Architectures}

\subsection{Memory Architectures}

\mult{2}

\begin{definition}{Three Architectures}\\
    \begin{itemize}
        \item \textbf{Von Neumann}: shared I+D bus
        \item \textbf{Harvard}: separate I/D
        \item \textbf{Modified Harvard}: split L1, unified L2/RAM
    \end{itemize}
\end{definition}

\begin{concept}{Parallelising a CPU}\\
    \begin{itemize}
        \item multicore
        \item multiple FUs (data parallelism)
        \item multiple pipelines (instruction)
        \item heterogeneous PEs
    \end{itemize}
\end{concept}

\multend

\begin{KR}{Recipe: Classify a CPU Architecture}\\
    shared bus $\to$ Von Neumann; fully separate $\to$ Harvard; split L1 + unified L2 $\to$ \important{Modified Harvard} (e.g.\ A53).
\end{KR}

\begin{example2}{Ex.\ S4--7 -- A53 architecture}\\
    Separate L1-I/L1-D, unified L2 -- which architecture? Ways to parallelise?
    \tcblower
    \important{Modified Harvard} (Harvard inside, Von Neumann outside). Parallelise via multicore, FUs, pipelines, PEs.
\end{example2}

\subsection{Multiple Functional Units (VLIW)}

\begin{definition}{VLIW}\\
    The \emph{compiler} packs one op per FU (.M/.D/.L/.S) into a wide instruction word, run in parallel. Big code, binary-incompatible; + software pipelining.
\end{definition}

\begin{example2}{Ex.\ S8--10 -- DSP}\\
    What memory architecture? How does execution work?
    \tcblower
    \important{Harvard} (separate program/data bus keeps MAC fed); \important{VLIW} execution + software pipelining.
\end{example2}

\subsection{Pipelines}

\mult{2}

\begin{definition}{Hazards}\\
    \begin{itemize}
        \item structural, data, control
        \item RISC = pipeline-friendly; CISC = microcoded
    \end{itemize}
\end{definition}

\begin{concept}{Pipeline \& Issue}\\
    \begin{itemize}
        \item super-pipeline: finer stages (faster)
        \item superscalar: more pipelines (several/cycle)
        \item order: in-order $\to$ OoO exec $\to$ OoO issue+exec
    \end{itemize}
\end{concept}

\multend

\begin{KR}{Recipe: Identify the Parallelism Mechanism}\\
    multiple FUs (VLIW)=data; multiple pipelines (superscalar)=instruction; idle slots filled by 2nd thread=SMT; separate cores=multicore/AMP.
\end{KR}

\begin{example2}{Ex.\ S11--14 -- Pipelines}\\
    Super-pipeline vs superscalar? Issue/execute order?
    \tcblower
    Super-pipeline = finer stages (issue faster); superscalar = more pipelines (several/cycle). Order: in-order $\to$ OoO exec $\to$ \important{OoO issue+exec}.
\end{example2}

\subsection{Co-processors \& SMT}

\begin{definition}{Co-processor}\\
    CPU passes unrecognised instructions to it and \emph{stalls} (8087 + 8086 on \texttt{FMUL}); own ISA/clock/mem (NEON). Costly decoders.
\end{definition}

\begin{concept}{SMT (Hyperthreading)}\\
    2 threads/core fill \important{idle issue slots}. Requires a \emph{superscalar} core (spare slots). $\sim$30\% gain; security issues.
\end{concept}

\begin{example2}{Ex.\ S15--19 -- Co-processor \& why SMT needs superscalar}\\
    How does a co-processor work? Why does SMT need superscalar?
    \tcblower
    Co-pro: CPU offloads + stalls. SMT fills idle slots with a 2nd thread; only superscalar cores have multiple slots/cycle to fill -- a scalar core has nothing spare.
\end{example2}

\subsection{Asymmetric Multiprocessing}

\begin{concept}{AMP Scenarios}\\
    different/same ISA, different clocks, big.LITTLE (same ISA), different/no OS.
\end{concept}

\begin{definition}{openAMP}\\
    \begin{itemize}
        \item \texttt{remoteproc}: boots a remote core's ELF
        \item \texttt{rpmsg} over \texttt{virtio}: ring-buffer messaging
    \end{itemize}
\end{definition}

\begin{KR}{Recipe: Set Up AMP}\\
    same ISA $\to$ SMP/big.LITTLE one OS; different ISA $\to$ separate OS/bare-metal per core + openAMP (\texttt{remoteproc}, \texttt{rpmsg}/\texttt{virtio}).
\end{KR}

\begin{example2}{Ex.\ S20--25 -- Different ISA under one OS?}\\
    Can different-ISA cores share one OS? How do they communicate?
    \tcblower
    \important{No} (one kernel = one ISA; big.LITTLE works only because same ISA). Different ISA $\to$ separate OSes + openAMP. Zynq: 4$\times$A53 SMP + 2$\times$R5 lockstep + FPGA.
\end{example2}
